% !TEX root = ../../main.tex
% C9.3 — Hướng phát triển. Có căn cứ (report-data-guide.md C9): xếp hạng lại image–text của RaSa, tinh chỉnh cùng miền, OpenVINO (R10), đánh giá BoT-SORT;
% vector CLIP thứ hai mỗi lần xuất hiện (C8.4, B4); bộ mã hóa dùng chung. Đồng bộ với bản tiếng Anh 12.3. Các hướng mở rộng khác lấy từ phạm vi Mục 1.3 (ngoài phạm vi).
\section{Hướng phát triển trong tương lai}
\label{sec:9.3}

Hướng phát triển ưu tiên là cải thiện tìm kiếm bằng văn bản và thuộc tính. Bước xếp hạng lại bằng bộ mã hóa đa phương thức của RaSa nay là một bước tùy chọn của luồng tìm kiếm và là nơi chứa độ chính xác của mô hình (Mục~\ref{sec:8.4}); để nó đủ nhanh cho sử dụng hằng ngày cần tính token ảnh của mỗi lần xuất hiện ngay lúc lập chỉ mục, trong tiến trình xử lý nền, thay vì ở lần tìm đầu tiên, và chạy các lượt hợp nhất trên GPU, nơi 128 ứng viên mất chưa tới một giây. Phép đo ở Mục~\ref{sec:8.4} cũng chỉ ra một hướng đơn giản hơn: một vector thứ hai cho mỗi lần xuất hiện từ bộ mã hóa kiểu CLIP, tính lúc lập chỉ mục và lưu trong collection Milvus riêng, sẽ cho tìm kiếm bằng văn bản và thuộc tính chất lượng gần với khi xếp hạng lại mà chỉ tốn thêm một lượt tìm kiếm vector, trong khi vector của RaSa tiếp tục phục vụ truy vấn bằng ảnh; thiết kế vốn đã giữ mỗi phiên bản bộ mã hóa một collection, nên thay đổi gói gọn trong quy trình lập chỉ mục và bộ định tuyến truy vấn. Tinh chỉnh một trong hai mô hình trên dữ liệu cùng miền với camera giám sát là bước sau đó.

Về bộ nhớ, một tiến trình bộ mã hóa dùng chung cho máy chủ ứng dụng và tiến trình xử lý nền sẽ bỏ được một trong hai bản bộ mã hóa truy vấn, khoảng 1~GiB. Về tốc độ, suy luận trên CPU có thể được tối ưu, chẳng hạn bằng OpenVINO, trước hết cho bộ mã hóa ảnh vốn chiếm phần lớn thời gian xử lý (Mục~\ref{sec:8.5}), hoặc hệ thống có thể được triển khai trên máy có GPU để phân tích kịp nhiều luồng camera cùng lúc thay vì luân phiên. Khi đã theo kịp thời gian thực, mỗi camera có thể được xử lý liên tục, và mỗi lần xuất hiện được mã hóa, ghi ngay khi kết thúc để rút ngắn độ trễ tới lúc tìm kiếm được. Bước ghi vào ba kho vốn đã xử lý riêng từng lần xuất hiện, nên thay đổi chủ yếu nằm ở quy trình phân tích.

Phần đánh giá cũng nên so sánh ByteTrack với BoT-SORT. Xa hơn, hệ thống có thể được mở rộng để liên kết các lần xuất hiện của cùng một người giữa nhiều camera, biểu diễn mỗi lần xuất hiện bằng nhiều khung hình, và hỗ trợ truy vấn bằng tiếng Việt.
